% Obsahový zdroj: import_hodiny/39_Binomicka_veta.tex a odpovídající PDF (obě strany).
% P1–P5: zdrojové Příklady 1–5; C1–C5: zdrojové Procvičování 1–5.
% Matematický obsah zachován; výklad, odvození a řešení pedagogicky rozvedeny.
% Upřesnění a opravy: viz KMB_29-41_Kontrola.md; původní import zůstává beze změny.
% Rozpis: 4_rocnik/tematicke_plany/Rozpis_vyuky_MAT_4_rocnik.tex, hodina 39.
% Plán: 4_rocnik/tematicke_plany/SPS.pdf, oddíly 2.1–2.3.
% Sazba: sablony/pripravy/TEMPLATE_Priprava.tex a společný priprava.sty.
% Vlastní doplnění: didaktický výklad, zdůvodnění, metodika; žádné nové číslované úlohy.
% Učebnice/09_Kombinatorika nemá obsahovou kapitolu; místní kombinatorické listy ani Realisticky.cz nejsou dostupné.
\documentclass[11pt,a4paper]{article}
% Spolecny styl pro ucitelske pripravy MAT (Work/MAT)
% Verze 1.1, 2026-10-23
\NeedsTeXFormat{LaTeX2e}
\ProvidesPackage{priprava}[2026/10/23 v1.1 Sablona ucitelskych priprav MAT]

% Zakladni typografie a cestina
\RequirePackage[T1]{fontenc}
\RequirePackage[utf8]{inputenc}
\RequirePackage[czech]{babel}
\RequirePackage{lmodern}
\RequirePackage{microtype}
\RequirePackage{amsmath,amssymb}
\RequirePackage{xcolor}
\RequirePackage{enumitem}
\RequirePackage[most]{tcolorbox}
\RequirePackage{tabularx,booktabs,array}
\RequirePackage{fancyhdr}
\RequirePackage{needspace}
\RequirePackage{geometry}
\RequirePackage[hidelinks]{hyperref}

\geometry{a4paper,top=17mm,bottom=18mm,left=23mm,right=23mm}
\definecolor{notegray}{gray}{0.39}
\setlength{\parindent}{0pt}
\setlength{\parskip}{2.5pt plus 1pt minus 0.5pt}
\setlist[itemize]{topsep=2pt,itemsep=1pt,parsep=0pt,leftmargin=1.45em}
\setlist[enumerate]{topsep=2pt,itemsep=2pt,parsep=0pt,leftmargin=1.65em}
\setcounter{secnumdepth}{0}
\renewcommand{\arraystretch}{1.12}

% Zahlavi: vlevo rocnik | globalni tema, vpravo cislo: nazev hodiny
% Zapati zustava zcela prazdne.
\pagestyle{fancy}
\fancyhf{}
\renewcommand{\headrulewidth}{0.35pt}
\renewcommand{\footrulewidth}{0pt}
\setlength{\headheight}{14pt}
\newcommand{\nastavPripravu}[4]{%
  \fancyhead[L]{\small #1\enspace\textbar\enspace #2}%
  \fancyhead[R]{\small #3: #4}%
}

% Minimalisticke cernobile boxy. Kazda definice zvlast.
\tcbset{commonlesson/.style={enhanced,breakable=false,colback=white,colframe=black,
  colbacktitle=white,coltitle=black,boxrule=0.45pt,
  arc=0.8mm,outer arc=0.8mm,left=3mm,right=3mm,top=1.6mm,bottom=1.6mm,
  toptitle=1mm,bottomtitle=0.7mm,fonttitle=\bfseries,
  attach boxed title to top left={xshift=2.5mm,yshift=-1.15mm},
  boxed title style={colback=white,colframe=white,boxrule=0pt,left=1mm,right=1mm}}}
\newtcolorbox{zadaniBox}[1][]{commonlesson,title={Zadání},#1}
\newtcolorbox{prikladBox}[1][]{commonlesson,title={Příklad},#1}
\newtcolorbox{definiceBox}[1][]{commonlesson,title={Definice},#1}

% Spolecny box pro zadani a jeho reseni, bez mezery mezi castmi.
% Pouziti:
%   \begin{prikladReseniBox}
%     Zadani ulohy.
%     \tcblower
%     \begin{align*} ... \end{align*}
%   \end{prikladReseniBox}
% Spodni cast ma jemny sedy podklad a delici linku; lze v ni sazet
% vice radku matematiky. Obe casti maji jen jeden spolecny ramecek.
\tcbset{solutionlesson/.style={
  commonlesson,
  enhanced jigsaw,
  breakable,
  colbacklower=black!3,
  segmentation style={solid,draw=black!25,line width=0.35pt},
  fontlower=\small\color{notegray},
  before lower={\noindent\textbf{Řešení:}\enspace},
  leftlower=3mm,rightlower=3mm
}}
\newtcolorbox{prikladReseniBox}[1][]{solutionlesson,title={Příklad},#1}
\newtcolorbox{zadaniReseniBox}[1][]{solutionlesson,title={Zadání},#1}

% Pomucky pro autora pripravy: sedive poznamky; \reseni pro starsi pripravy
\newcommand{\poznamka}[1]{\par\begingroup\small\color{notegray}\itshape
  \textnormal{\textbf{Poznámka pro učitele.}}\enspace #1\par\endgroup}
\newcommand{\reseni}[1]{\par\begingroup\small\color{notegray}
  \textbf{Očekávané řešení:}\enspace #1\par\endgroup}
\newcommand{\blok}[1]{\needspace{5\baselineskip}\subsection*{#1}}
\newcommand{\vyrok}[1]{\uv{#1}}

\nastavPripravu{4. ročník}{Kombinatorika}{39}{Binomická věta}
\begin{document}


Připomeň $(a+b)^2$ a $(a+b)^3$ a porovnej jejich koeficienty s řádky Pascalova trojúhelníku. Cílem je vysvětlit tento jev výběrem členů při násobení závorek.
\poznamka{Předpoklady: násobení mnohočlenů, mocniny součinu, kombinační čísla a Pascalův trojúhelník. Cíl: žák zdůvodní koeficienty, rozvine malou mocninu a bez celého rozvoje najde zadaný koeficient. Pro společnou práci upřednostni P1–P3; P4 je náročnější pokračování, P5 krátké propojení s hodinou 37.}

\blok{Výklad a odvození}

\textbf{Binomická věta.} Pro reálná čísla $a,b$ a celé $n\geq1$ platí
\[
 (a+b)^n=\sum_{k=0}^n\binom nk a^{n-k}b^k.
\]
Věta platí i jako rovnost polynomů a pro $n=0$ má význam rovnosti konstantních polynomů $1=1$. V zápisu polynomu nahrazujeme činitel s exponentem nula konstantou $1$; nevzniká tím úloha na číselný výraz $0^0$.

\textbf{Odvození.} Rozepišme mocninu na součin $n$ závorek $(a+b)$. V každé při roznásobení vybereme buď $a$, nebo $b$. Aby vznikl součin $a^{n-k}b^k$, musíme zvolit $b$ právě v $k$ z těchto $n$ závorek. Výběr pozic lze provést $\binom nk$ způsoby. Rozdělení podle počtu vybraných $b$ pokryje všechna roznásobení, proto koeficienty sčítaných členů mají právě uvedený tvar.

Člen odpovídající indexu $k$ je $\binom nk a^{n-k}b^k$ a stojí na pozici $k+1$. Exponent prvního členu klesá od $n$ k nule, exponent druhého roste od nuly k $n$; jejich součet je stále $n$. Koeficienty čteme z řádku s indexem $n$. Symbolický rozvoj v nezávislých proměnných $a,b$ obsahuje $n+1$ takových členů; po konkrétním dosazení se některé mohou sloučit nebo vynulovat.\begin{definiceBox}
\textbf{Koeficient u mocniny $x^r$} v mnohočlenu je číslo, kterým je tato mocnina násobena po sloučení stejných mocnin.
\end{definiceBox}
\begin{definiceBox}
\textbf{Konstantní člen} je člen nezávislý na $x$, tedy koeficient u $x^0$.
\end{definiceBox}

Pro rozvoj rozdílu zapisujeme celý druhý člen se znaménkem: $(a-c)^n=(a+(-c))^n$. V obecném členu se pak umocňuje $(-c)^k$. Obdobně $(2x)^{n-k}$ znamená $2^{n-k}x^{n-k}$; číselný faktor nelze z mocniny vynechat.

Při hledání koeficientu nejprve napiš obecný člen, potom pomocí požadované mocniny najdi $k$, až pak vypočti číselný koeficient. Index $k$ nemusí být roven hledanému exponentu $x$: záleží na tom, v kterém členu dvojčlenu proměnná stojí.


\poznamka{Průběh: 0–11 min odvození; 11–24 min P1–P2; 24–34 min P3 a podle tempa P4; 34–42 min C1–C3; 42–45 min kontrola a krátké P5. C4–C5 a neprobraný P4 jsou rezerva či domácí práce. Znaménka a mocniny koeficientů mají přednost před množstvím hotových rozvojů.}

\blok{Společně řešené příklady}

% Zdroj P1: import_hodiny/39_Binomicka_veta.tex, Příklad 1.
\begin{prikladReseniBox}[title={Příklad P1}]
Rozviň \((x+2)^4\).
\tcblower
Použijeme $a=x$, $b=2$, $n=4$. Koeficienty jsou $1,4,6,4,1$:
\begin{align*}
 (x+2)^4&=x^4+4x^3\cdot2+6x^2\cdot2^2+4x\cdot2^3+2^4\\
 &=x^4+8x^3+24x^2+32x+16.
\end{align*}
Kontrola pro $x=1$: levá strana $3^4=81$, pravá $1+8+24+32+16=81$.
\end{prikladReseniBox}


% Zdroj P2: import_hodiny/39_Binomicka_veta.tex, Příklad 2.
\begin{prikladReseniBox}[title={Příklad P2}]
Rozviň \((2x-y)^3\).
\tcblower
Dosadíme $a=2x$, $b=-y$, $n=3$. Umocňujeme celé oba členy:
\begin{align*}
 (2x-y)^3&=(2x)^3+3(2x)^2(-y)+3(2x)(-y)^2+(-y)^3\\
 &=8x^3-12x^2y+6xy^2-y^3.
\end{align*}
Znaménka určuje parita exponentu u $-y$. Pro $x=y=1$ vychází $1=8-12+6-1$.
\end{prikladReseniBox}


% Zdroj P3: import_hodiny/39_Binomicka_veta.tex, Příklad 3.
\begin{prikladReseniBox}[title={Příklad P3}]
Najdi koeficient u \(x^3\) v \((x+2)^5\).
\tcblower
Obecný člen je $\binom5k x^{5-k}2^k$. Požadujeme $5-k=3$, tedy $k=2$. Člen je $\binom52\,2^2x^3=40x^3$, takže hledaný koeficient je $40$. Číslo $10$ z Pascalova trojúhelníku je jen jedna část koeficientu.
\end{prikladReseniBox}


% Zdroj P4: import_hodiny/39_Binomicka_veta.tex, Příklad 4.
\begin{prikladReseniBox}[title={Příklad P4}]
Najdi koeficient u \(x^4\) v \((2x-1)^6\).
\tcblower
Obecný člen má tvar
\[
 \binom6k(2x)^{6-k}(-1)^k
 =\binom6k2^{6-k}(-1)^k x^{6-k}.
\]
Pro mocninu $x^4$ musí být $6-k=4$, tedy $k=2$. Hledaný koeficient je $\binom62\,2^4(-1)^2=15\cdot16=240$. Znaménko je kladné, protože $k$ je sudé.
\end{prikladReseniBox}


% Zdroj P5: import_hodiny/39_Binomicka_veta.tex, Příklad 5.
\begin{prikladReseniBox}[title={Příklad P5}]
Použij binomickou větu ke zdůvodnění \(\sum_{k=0}^n\binom nk=2^n\).
\tcblower
Dosadíme $a=b=1$. Potom
\[
 2^n=(1+1)^n=\sum_{k=0}^n\binom nk1^{n-k}1^k
 =\sum_{k=0}^n\binom nk.
\]
Pro $n=0$ dostáváme $1=1$. Jde o druhé zdůvodnění identity, kterou jsme v hodině 37 odvodili počítáním všech podmnožin.
\end{prikladReseniBox}


\newpage
\blok{Samostatná práce a zásoba k procvičení}

U každé úlohy uveď substituci za $a,b$, řádek koeficientů nebo obecný člen. C4–C5 využij pro rychlejší studenty nebo jako návazné procvičení.

% Zdroj C1: import_hodiny/39_Binomicka_veta.tex, Procvičování, položka 1.
\begin{zadaniReseniBox}[title={Úloha C1}]
Rozviň \((a+b)^3\).
\tcblower
Pro $n=3$ jsou koeficienty $1,3,3,1$; mocnina $a$ klesá a mocnina $b$ roste:
\[
(a+b)^3=a^3+3a^2b+3ab^2+b^3.
\]
\end{zadaniReseniBox}

% Zdroj C2: import_hodiny/39_Binomicka_veta.tex, Procvičování, položka 2.
\begin{zadaniReseniBox}[title={Úloha C2}]
Rozviň \((x-1)^4\).
\tcblower
Dosadíme $a=x$, $b=-1$, $n=4$. Proto
\begin{align*}
(x-1)^4&=x^4+4x^3(-1)+6x^2(-1)^2+4x(-1)^3+(-1)^4\\
&=x^4-4x^3+6x^2-4x+1.
\end{align*}
\end{zadaniReseniBox}

% Zdroj C3: import_hodiny/39_Binomicka_veta.tex, Procvičování, položka 3.
\begin{zadaniReseniBox}[title={Úloha C3}]
Najdi koeficient u \(x^2\) v \((x+3)^4\).
\tcblower
Obecný člen je $\binom4k x^{4-k}3^k$. Z $4-k=2$ vychází $k=2$, takže koeficient je $\binom42\,3^2=6\cdot9=54$.
\end{zadaniReseniBox}

% Zdroj C4: import_hodiny/39_Binomicka_veta.tex, Procvičování, položka 4.
\begin{zadaniReseniBox}[title={Úloha C4}]
Najdi konstantní člen v \((2x+5)^3\).
\tcblower
Konstantní člen vznikne volbou $5$ ve všech třech závorkách, tedy má hodnotu $5^3=125$. Také dosazení $x=0$ do polynomu dává jeho konstantní člen: $(2\cdot0+5)^3=125$.
\end{zadaniReseniBox}

% Zdroj C5: import_hodiny/39_Binomicka_veta.tex, Procvičování, položka 5.
\begin{zadaniReseniBox}[title={Úloha C5}]
Najdi koeficient u \(x^3\) v \((1-2x)^5\).
\tcblower
Zvolíme $a=1$, $b=-2x$, $n=5$. Obecný člen je $\binom5k1^{5-k}(-2x)^k$. Pro $x^3$ potřebujeme $k=3$, koeficient je $\binom53(-2)^3=10\cdot(-8)=-80$.
\end{zadaniReseniBox}


\poznamka{Typické chyby: $(a+b)^n=a^n+b^n$; chybějící mocnina dvojky v $(2x)^r$; záměna indexu $k$ a pořadí členu; znaménko $(-2)^3$. Dosazení konkrétního $x$ slouží ke kontrole, ale samo nedokazuje rovnost celých polynomů. Hodina 40 ověří variace, permutace a kombinace; binomická věta není povinnou součástí tohoto testu. V hodině 41 se k ní lze krátce vrátit v úloze B1.}
\end{document}
